% =========================================================
% 11_change_of_basis.tex
% Vectors and Linear Algebra
% =========================================================

% =========================================================
% SECTION DIVIDER
% =========================================================

\section{Change of Basis}

% =========================================================
% BASIS AND COORDINATES
% =========================================================

\begin{frame}{Basis and Coordinates}

A vector can be represented using different bases.

Consider the standard basis of $\mathbb{R}^2$:
\vspace{-0.1cm}
\[
\mathcal{E}
=
\left\{
\begin{bmatrix}
1\\0
\end{bmatrix},
\begin{bmatrix}
0\\1
\end{bmatrix}
\right\}.
\]

For
\vspace{-0.1cm}
\[
\vect{v}
=
\begin{bmatrix}
3\\2
\end{bmatrix},
\]

its coordinates in the standard basis are
\vspace{-0.1cm}
\[
[\vect{v}]_{\mathcal{E}}
=
\begin{bmatrix}
3\\2
\end{bmatrix}.
\]

The vector itself and its coordinate representation
are related, but are not the same concept.

\end{frame}


% =========================================================
% NON-STANDARD BASIS
% =========================================================

\begin{frame}{A Different Basis}
\vspace{-0.3cm}
Consider the basis
\vspace{-0.1cm}
\[
\mathcal{B}
=
\left\{
\begin{bmatrix}
1\\1
\end{bmatrix},
\begin{bmatrix}
1\\-1
\end{bmatrix}
\right\}.
\]

We want to express
\vspace{-0.4cm}
\[
\vect{v}
=
\begin{bmatrix}
3\\1
\end{bmatrix}
\]
\vspace{-0.5cm}

using this basis.

Write
\vspace{-0.1cm}
\[
\vect{v}
=
c_1
\begin{bmatrix}
1\\1
\end{bmatrix}
+
c_2
\begin{bmatrix}
1\\-1
\end{bmatrix}.
\]
\vspace{-0.5cm}

This gives
\vspace{-0.1cm}
\[
c_1+c_2=3,
\qquad
c_1-c_2=1.
\]

Hence,
\vspace{-0.1cm}
\[
c_1=2,
\qquad
c_2=1.
\]

\vspace{-0.3cm}

Therefore,
\vspace{-0.1cm}
\[
\boxed{
[\vect{v}]_{\mathcal{B}}
=
\begin{bmatrix}
2\\1
\end{bmatrix}.
}
\]

\end{frame}


% =========================================================
% VECTOR VS COORDINATES
% =========================================================

\begin{frame}{Vector vs. Coordinate Representation}

For the same vector $\vect{v}$,

\[
[\vect{v}]_{\mathcal{E}}
=
\begin{bmatrix}
3\\1
\end{bmatrix}
\]

while

\[
[\vect{v}]_{\mathcal{B}}
=
\begin{bmatrix}
2\\1
\end{bmatrix}.
\]

The numerical coordinates are different because the bases
are different.

\vspace{0.4cm}

However,

\[
\boxed{
\text{The vector itself has not changed.}
}
\]

Only its representation has changed.

\end{frame}


% =========================================================
% BASIS AND COORDINATES
% =========================================================

\begin{frame}{Basis and Coordinates}

Let

\[
\mathcal{B}
=
\{\vect{b}_1,\vect{b}_2,\ldots,\vect{b}_n\}
\]

be a basis for $\mathbb{R}^n$.

\vspace{0.1cm}

Every vector $\vect{v}\in\mathbb{R}^n$ can be written uniquely as
\vspace{-0.1cm}
\[
\boxed{
\vect{v}
=
c_1\vect{b}_1
+
c_2\vect{b}_2
+
\cdots
+
c_n\vect{b}_n
}
\]

\vspace{0.1cm}

The coefficients form the coordinate vector of $\vect{v}$
with respect to $\mathcal{B}$:
\vspace{-0.1cm}
\[
\boxed{
[\vect{v}]_{\mathcal{B}}
=
\begin{bmatrix}
c_1\\
c_2\\
\vdots\\
c_n
\end{bmatrix}
}
\]

\end{frame}

% =========================================================
% BASIS MATRIX
% =========================================================

\begin{frame}{Basis Matrix}

The basis vectors can be arranged as the columns of a matrix:

\[
\boxed{
P_{\mathcal{B}}
=
\begin{bmatrix}
|&|&&|\\
\vect{b}_1&\vect{b}_2&\cdots&\vect{b}_n\\
|&|&&|
\end{bmatrix}
}
\]

\vspace{0.3cm}

Then the vector can be reconstructed from its coordinates:

\[
\boxed{
\vect{v}
=
P_{\mathcal{B}}[\vect{v}]_{\mathcal{B}}
}
\]

\vspace{0.3cm}

\begin{block}{Interpretation}
$P_{\mathcal{B}}$ converts coordinates expressed in basis
$\mathcal{B}$ into the standard-coordinate representation.
\end{block}

\end{frame}


% =========================================================
% CHANGE TO STANDARD COORDINATES
% =========================================================

\begin{frame}{From Basis Coordinates to Standard Coordinates}

Consider the basis
\vspace{-0.1cm}
\[
\mathcal{B}
=
\left\{
\begin{bmatrix}
1\\
1
\end{bmatrix},
\begin{bmatrix}
1\\
-1
\end{bmatrix}
\right\}.
\]

Its basis matrix is
\vspace{-0.1cm}
\[
\boxed{
P_{\mathcal{B}}
=
\begin{bmatrix}
1&1\\
1&-1
\end{bmatrix}
}
\]

\vspace{0.1cm}

Suppose the coordinates of $\vect{v}$ in $\mathcal{B}$ are
\vspace{-0.1cm}
\[
[\vect{v}]_{\mathcal{B}}
=
\begin{bmatrix}
2\\
1
\end{bmatrix}.
\]

Then
\vspace{-0.1cm}
\[
\vect{v}
=
P_{\mathcal{B}}[\vect{v}]_{\mathcal{B}}
=
\begin{bmatrix}
1&1\\
1&-1
\end{bmatrix}
\begin{bmatrix}
2\\
1
\end{bmatrix}
=
\boxed{
\begin{bmatrix}
3\\
1
\end{bmatrix}
}
\]

\end{frame}


% =========================================================
% INVERSE BASIS MATRIX
% =========================================================

\begin{frame}{From Standard Coordinates to Basis Coordinates}

Starting from
\vspace{-0.1cm}
\[
\vect{v}
=
P_{\mathcal{B}}[\vect{v}]_{\mathcal{B}},
\]

multiply by $P_{\mathcal{B}}^{-1}$:
\vspace{-0.1cm}
\[
P_{\mathcal{B}}^{-1}\vect{v}
=
[\vect{v}]_{\mathcal{B}}.
\]

Therefore,
\vspace{-0.1cm}
\[
\boxed{
[\vect{v}]_{\mathcal{B}}
=
P_{\mathcal{B}}^{-1}\vect{v}.
}
\]

Thus,
\vspace{-0.1cm}
\[
\boxed{
\begin{array}{ccc}
[\vect{v}]_{\mathcal{B}}
&\xrightarrow{\quad P_{\mathcal{B}}\quad}&
\vect{v}\\[4pt]
\vect{v}
&\xrightarrow{\quad P_{\mathcal{B}}^{-1}\quad}&
[\vect{v}]_{\mathcal{B}}
\end{array}
}
\]

\end{frame}


% =========================================================
% CHANGE OF BASIS
% =========================================================

\begin{frame}{Change of Basis}

Suppose $\mathcal{B}$ and $\mathcal{C}$ are two bases.

We can convert coordinates from $\mathcal{B}$ to $\mathcal{C}$.

Since
\vspace{-0.1cm}
\[
\vect{v}
=
P_{\mathcal{B}}[\vect{v}]_{\mathcal{B}}
\]

and

\[
[\vect{v}]_{\mathcal{C}}
=
P_{\mathcal{C}}^{-1}\vect{v},
\]

we obtain

\[
\boxed{
[\vect{v}]_{\mathcal{C}}
=
P_{\mathcal{C}}^{-1}
P_{\mathcal{B}}
[\vect{v}]_{\mathcal{B}}.
}
\]

Therefore, the change-of-basis matrix is

\[
\boxed{
P_{\mathcal{C}\leftarrow\mathcal{B}}
=
P_{\mathcal{C}}^{-1}P_{\mathcal{B}}.
}
\]

\end{frame}


% =========================================================
% CHANGE OF BASIS EXAMPLE
% =========================================================

\begin{frame}{Example: Changing Coordinates}

Let
\vspace{-0.2cm}
\[
\mathcal{B}
=
\left\{
\begin{bmatrix}
1\\1
\end{bmatrix},
\begin{bmatrix}
1\\-1
\end{bmatrix}
\right\}
\]

and let $\mathcal{E}$ be the standard basis.

Then
\vspace{-0.4cm}
\[
P_{\mathcal{B}}
=
\begin{bmatrix}
1&1\\
1&-1
\end{bmatrix}.
\]
\vspace{-0.4cm}

Suppose
\vspace{-0.1cm}
\[
[\vect{v}]_{\mathcal{B}}
=
\begin{bmatrix}
2\\1
\end{bmatrix}.
\]
\vspace{-0.4cm}

Then
\vspace{-0.1cm}
\[
\vect{v}
=
P_{\mathcal{B}}[\vect{v}]_{\mathcal{B}}
=
\begin{bmatrix}
3\\1
\end{bmatrix}.
\]
\vspace{-0.4cm}

Hence,
\vspace{-0.1cm}
\[
\boxed{
[\vect{v}]_{\mathcal{E}}
=
\begin{bmatrix}
3\\1
\end{bmatrix}.
}
\]

\end{frame}


% =========================================================
% TRANSFORMATION IN DIFFERENT BASES
% =========================================================

\begin{frame}{Linear Transformations in Different Bases}

Suppose
\vspace{-0.1cm}
\[
T:V\rightarrow V
\]

is a linear transformation. Its matrix depends on the chosen basis.

Let
\vspace{-0.1cm}
\[
A_{\mathcal{B}}
\]

represent $T$ in basis $\mathcal{B}$, and
\vspace{-0.1cm}
\[
A_{\mathcal{C}}
\]

represent the same transformation in basis $\mathcal{C}$.

Then
\vspace{-0.1cm}
\[
\boxed{
A_{\mathcal{C}}
=
P_{\mathcal{C}}^{-1}
A_{\mathcal{B}}
P_{\mathcal{C}}.
}
\]

Thus, changing the basis changes the matrix representation,
not the underlying transformation.

\end{frame}


% =========================================================
% SIMILAR MATRICES
% =========================================================

\begin{frame}{Similar Matrices}

Two square matrices $A$ and $B$ are \textbf{similar} if

\[
\boxed{
B=P^{-1}AP
}
\]

for some invertible matrix $P$.

\vspace{0.4cm}

Similar matrices represent the same linear transformation
under different coordinate systems.

\vspace{0.4cm}

Important properties are preserved under similarity,
including

\[
\boxed{
\det(A)=\det(B).
}
\]

The eigenvalues are also preserved.

\end{frame}


% =========================================================
% WHY CHANGE BASIS?
% =========================================================

\begin{frame}{Why Change the Basis?}

A suitable basis can make a linear transformation easier
to understand or compute.

For example, a matrix may have a complicated representation
in the standard basis but become diagonal in another basis.

\vspace{0.4cm}

This motivates the idea of

\[
\boxed{
\text{choosing coordinates that simplify the problem}.
}
\]

\vspace{0.4cm}

This idea leads directly to

\[
\boxed{
\text{eigenvectors}
\quad\text{and}\quad
\text{diagonalization}.
}
\]

\end{frame}